\chapter{المدارات الحدودية والتفاعلية}\label{ch:b2:frontier-orbitals}

قارورة من البنتاديين الحلقي تُترك على الرف تتحول ببطء إلى ثنائيها:
يتحد جزيئان من \omterm{def:b2:frontier-orbitals:diels-alder}{الديين} نفسه في مركب ثنائي الحلقة ذي جسر، 
ولا يتكوّن إلا واحد من المتماكبات الفراغية الممكنة. ولا شيء في الأسهم المنحنية في
مجلد السنة الأولى يتنبأ بأيها. أما مداران فيتنبآن به: أعلى مدار مشغول
في أحد الجزيئين وأدنى مدار فارغ في الآخر. ويحوّل هذا
الفصل تآثر المستويين في \cref{ch:b2:diatomic-mos} إلى
أداة لدراسة التفاعلية --- أي ذرة يهاجمها المحب للنواة، ومن أي جهة، 
ولماذا يتحد \omterm{def:b2:frontier-orbitals:diels-alder}{ديين} وألكين في حلقة في خطوة واحدة.

\begin{recall}
\Cref{ch:b2:diatomic-mos}: مداران متآثران، وتنافر الإلكترونات الأربعة.
و\cref{ch:b2:huckel}: مستويات هوكل ومعاملاتها في الأنظمة المترافقة.
ومجلد السنة الأولى: المحبات للنواة والمحبات للإلكترونات، والأسهم المنحنية، والتحكم
الحركي، والتفاعلات النوعية فراغيًا والانتقائية فراغيًا والانتقائية موضعيًا،
والإضافات المحبة للإلكترونات إلى الألكينات، والاستبدال المحب للنواة.
\end{recall}

\section{عودة إلى المدارين}

حين يقترب جزيئان، تتداخل مداراتهما وتتآثر. ومعظم
التآثر ضعيف، لأن المدارات متباعدة في الطاقة؛ فتكفي حينئذ صيغة اضطرابية من
نتيجة المستويين.

\begin{theorem}[التآثر الضعيف بين مستويين]\label{thm:b2:frontier-orbitals:perturbation}
مداران طاقتاهما $\varepsilon_1 < \varepsilon_2$، مقترنان \omterm{def:b2:diatomic-mos:integrals}{بتكامل رنين}
$\beta$ مع $|\beta| \ll \Delta\varepsilon = \varepsilon_2 - \varepsilon_1$ (مع إهمال التداخل)،
ينزاحان إلى
\[
  \varepsilon_1 - \frac{\beta^2}{\Delta\varepsilon}, \qquad \varepsilon_2 + \frac{\beta^2}{\Delta\varepsilon}
\]
حتى الرتبة الأولى في $\beta^2/\Delta\varepsilon^2$: ينخفض المستوى الأدنى ويرتفع الأعلى بالمقدار
نفسه، وهو أصغر كلما كانت الفجوة أكبر.
\end{theorem}

\begin{proof}
حسب \cref{thm:b2:diatomic-mos:heteronuclear} المستويات المضبوطة هي $\bar\varepsilon \mp
\sqrt{\Delta\varepsilon^2/4 + \beta^2}$. ومع $u = 4\beta^2/\Delta\varepsilon^2 \ll 1$، $\sqrt{\Delta\varepsilon^2/4
+ \beta^2} = \frac{\Delta\varepsilon}{2}\sqrt{1 + u} \approx \frac{\Delta\varepsilon}{2}(1 + u/2) =
\frac{\Delta\varepsilon}{2} + \frac{\beta^2}{\Delta\varepsilon}$، و$\bar\varepsilon \mp \Delta\varepsilon/2 = \varepsilon_1,
\varepsilon_2$.
\end{proof}

\begin{proposition}[إلكترونان يتجاذبان، وأربعة تتنافر]\label{prop:b2:frontier-orbitals:four-electron}
إذا حمل المداران المتآثران إلكترونين في المجموع، استقر النظام
بنحو $2\beta^2/\Delta\varepsilon$. وإذا حملا أربعة، فقد النظام استقراره.
\end{proposition}

\begin{proof}
يذهب إلكترونان إلى المستوى المنخفض: ربح قدره $2\beta^2/\Delta\varepsilon$. ومع أربعة، يمتلئ المستوى الأعلى
أيضًا؛ وحتى رتبة \cref{thm:b2:frontier-orbitals:perturbation} يتلاشى
الانزياحان، ومع الاحتفاظ بالتداخل يرتفع المستوى الأعلى أكثر مما ينخفض
الأدنى (\cref{prop:b2:diatomic-mos:asymmetry}): تنافر صافٍ.
\end{proof}

\begin{omfigure}
\begin{tikzpicture}[font=\small]
  \begin{scope}
    \pic (a) at (0,0) {omlevel={ud}{0.7}};
    \pic (b) at (3,1.4) {omlevel={empty}{0.7}};
    \pic (lo) at (1.5,-0.5) {omlevel={ud}{0.7}};
    \pic (hi) at (1.5,1.9) {omlevel={empty}{0.7}};
    \draw[omcorr, densely dashed] (a-r) -- (lo-l) (a-r) -- (hi-l) (b-l) -- (lo-r) (b-l) -- (hi-r);
    \node[below=6pt] at (1.5,-0.6) {\foreignlanguage{arabic}{إلكترونان: استقرار}};
    \node[left] at (a-l) {ممتلئ};
    \node[right] at (b-r) {فارغ};
  \end{scope}
  \begin{scope}[xshift=7cm]
    \pic (a) at (0,0) {omlevel={ud}{0.7}};
    \pic (b) at (3,0.6) {omlevel={ud}{0.7}};
    \pic (lo) at (1.5,-0.5) {omlevel={ud}{0.7}};
    \pic (hi) at (1.5,1.25) {omlevel={ud}{0.7}};
    \draw[omcorr, densely dashed] (a-r) -- (lo-l) (a-r) -- (hi-l) (b-l) -- (lo-r) (b-l) -- (hi-r);
    \node[below=6pt] at (1.5,-0.6) {\foreignlanguage{arabic}{أربعة إلكترونات: تنافر}};
    \node[left] at (a-l) {ممتلئ};
    \node[right] at (b-r) {ممتلئ};
  \end{scope}
\end{tikzpicture}

{\small تآثر مدارين لجزيئين متجاورين. المدار الممتلئ
المقابل لمدار فارغ يعطي استقرارًا صافيًا؛ والمداران الممتلئان يتنافران، إذ
يرتفع المستوى الأعلى أكثر مما ينخفض الأدنى. وهذا التنافر هو أصل
الإعاقة الفراغية بين الجزيئات.}
\end{omfigure}

\section{المدارات الحدودية}

\begin{definition}[المدارات الحدودية]\label{def:b2:frontier-orbitals:frontier}
المدار \emph{HOMO}\index{HOMO} (أعلى \omterm{def:b2:diatomic-mos:lcao}{مدار جزيئي} مشغول) والمدار
\emph{LUMO}\index{LUMO} (أدنى \omterm{def:b2:diatomic-mos:lcao}{مدار جزيئي} غير مشغول) لجزيء ما هما
\emph{مداراه الحدوديان}\index{مدارات حدودية}.
\end{definition}

\begin{proposition}[التآثر الغالب]\label{prop:b2:frontier-orbitals:dominant}
بين جزيئين ذوي طبقات مغلقة، تكون أهم التآثرات المثبِّتة
هي التآثرات بين \omterm{def:b2:frontier-orbitals:frontier}{HOMO} أحدهما \omterm{def:b2:frontier-orbitals:frontier}{وLUMO} الآخر؛ ومن بين الزوجين من هذا النوع، يغلب الزوج
ذو الفجوة الأصغر. فالمحب للنواة يتفاعل عبر مداره \omterm{def:b2:frontier-orbitals:frontier}{HOMO}، 
والمحب للإلكترونات عبر مداره \omterm{def:b2:frontier-orbitals:frontier}{LUMO}.
\end{proposition}

\begin{proof}[تعليل]
الأزواج ممتلئ--ممتلئ تتنافر (\cref{prop:b2:frontier-orbitals:four-electron}) والأزواج فارغ--فارغ
لا تحمل إلكترونات: فلا تثبّت إلا الأزواج ممتلئ--فارغ، كل منها بمقدار $2\beta^2/\Delta\varepsilon$.
والمدار \omterm{def:b2:frontier-orbitals:frontier}{HOMO} هو أعلى مستوى ممتلئ والمدار \omterm{def:b2:frontier-orbitals:frontier}{LUMO} أدنى مستوى فارغ، فيكون للزوج
HOMO--LUMO أصغر فجوة بينها؛ وحسب المبرهنة، تعطي أصغر فجوة
أكبر حد.
\end{proof}

\begin{definition}[التحكم بالشحنة والتحكم المداري]\label{def:b2:frontier-orbitals:control}
يكون التفاعل تحت \emph{التحكم بالشحنة}\index{تحكم بالشحنة} حين يقرر التجاذب
بين شحنات الشريكين موضع تفاعلهما، وتحت
\emph{التحكم المداري}\index{تحكم مداري} حين يقرره تداخل \omterm{def:b2:frontier-orbitals:frontier}{المدارات الحدودية}،
فيقع الهجوم حيث تكون معاملاتها أكبر.
\end{definition}

\begin{definition}[الأنواع القاسية واللينة]\label{def:b2:frontier-orbitals:hard-soft}
\emph{المحب للنواة القاسي}\index{محب للنواة قاس} أو \emph{المحب للإلكترونات القاسي}
\index{محب للإلكترونات قاس} صغير، يحمل شحنة مركزة وله
مدار \omterm{def:b2:frontier-orbitals:frontier}{HOMO} (أو \omterm{def:b2:frontier-orbitals:frontier}{LUMO} على التوالي) بعيد عن مدار شركائه: فيتفاعل تحت \omterm{def:b2:frontier-orbitals:control}{التحكم
بالشحنة}. و\emph{المحب للنواة اللين}\index{محب للنواة لين} أو \emph{المحب للإلكترونات اللين}
\index{محب للإلكترونات لين} كبير وقابل للاستقطاب، ذو مدار \omterm{def:b2:frontier-orbitals:frontier}{HOMO} مرتفع
(أو مدار \omterm{def:b2:frontier-orbitals:frontier}{LUMO} منخفض على التوالي): فيتفاعل تحت \omterm{def:b2:frontier-orbitals:control}{التحكم المداري}. ويتفاعل القاسي مفضِّلًا
القاسي، واللين اللين.
\end{definition}

أيون الهيدروكسيد وأيون الفلوريد محبان للنواة قاسيان، واليوديد والثيولات
لينة؛ والبروتون والكاتيون الكربوني محبان للإلكترونات قاسيان، وكربون
الألكان المهلجن لين.

\begin{definition}[المحب للنواة الثنائي الموقع]\label{def:b2:frontier-orbitals:ambident}
\emph{المحب للنواة الثنائي الموقع}\index{محب للنواة ثنائي الموقع} له ذرتان محبتان للنواة
مرتبطتان بالترافق، يمكن لكل منهما أن تهاجم محبًا للإلكترونات.
\end{definition}

يُهاجَم أيون السيانيد عند النيتروجين، حيث الشحنة السالبة أكبر، من قِبل \omterm{def:b2:frontier-orbitals:hard-soft}{المحبات للإلكترونات
القاسية}، وعند الكربون، حيث يكون معامل مداره \omterm{def:b2:frontier-orbitals:frontier}{HOMO} أكبر، من قِبل اللينة منها:
فيعطي الألكان المهلجن نتريلًا \ce{R-C#N}. وأيونات الإينولات، التي نلقاها في
\cref{ch:b2:enolates-aldol}، تتفاعل عند الكربون أو عند الأكسجين بالطريقة نفسها.

\begin{method}[تحليل بالمدارات الحدودية]\label{met:b2:frontier-orbitals:fmo}
\begin{enumerate}
  \item حدّد المحب للنواة (مدار \omterm{def:b2:frontier-orbitals:frontier}{HOMO} مرتفع) والمحب للإلكترونات (مدار \omterm{def:b2:frontier-orbitals:frontier}{LUMO} منخفض).
  \item ارسم المدارين الحدوديين بمعاملاتهما وإشاراتهما.
  \item يقع التفاعل حيث يكون التداخل أكبر: الذرات ذات المعاملات
        الأكبر، وفصوص متماثلة الإشارة متقابلة.
  \item يتبع اتجاه الاقتراب شكلَ المدار \omterm{def:b2:frontier-orbitals:frontier}{LUMO} (عموديًا على مستوي
        \ce{C=O}، ومن خلف الرابطة \ce{C-X}).
  \item كلما صغرت الفجوة HOMO--LUMO كان التفاعل أسرع (مع تساوي العوامل الأخرى).
\end{enumerate}
\end{method}

يُستنتج اتجاه الهجوم من المدار \omterm{def:b2:frontier-orbitals:frontier}{LUMO}. فالمدار \omterm{def:b2:frontier-orbitals:frontier}{LUMO} لمجموعة الكربونيل هو
مدارها $\pi^*$، الأكبر على الكربون والعمودي على المستوي: تقترب المحبات للنواة من
الكربون من فوق المستوي أو من تحته، لا على طول المحور \ce{C=O}. والمدار \omterm{def:b2:frontier-orbitals:frontier}{LUMO}
للألكان المهلجن هو المدار $\sigma^*$ للرابطة \ce{C-X}، الذي يتجه فصه الكبير على الكربون
بعيدًا عن X: فيهاجم المحب للنواة من الخلف، وينقلب الكربون
--- وهي الكيمياء الفراغية للاستبدال الثنائي الجزيئية في مجلد السنة الأولى.

\section{تفاعل ديلز--ألدر}

\begin{definition}[التفاعل المتضافر]\label{def:b2:frontier-orbitals:concerted}
\emph{التفاعل المتضافر}\index{تفاعل متضافر} يكوّن روابطه كلها ويكسرها في
خطوة واحدة، عبر حالة انتقالية واحدة، دون مركب وسيط.
\end{definition}

\begin{definition}[تفاعل ديلز--ألدر]\label{def:b2:frontier-orbitals:diels-alder}
\emph{تفاعل ديلز--ألدر}\index{تفاعل ديلز--ألدر} هو إضافة
\emph{ديين}\index{ديين} مترافق، في تشكّله s-cis، إضافةً متضافرة إلى ألكين أو
ألكاين، هو \emph{المحب للديين}\index{محب للديين}، فتتكوّن حلقة سداسية فيها رابطتان
$\sigma$ جديدتان ورابطة $\pi$ جديدة واحدة.
\end{definition}

\begin{omfigure}
\begin{tikzpicture}[font=\small, >={Stealth}, scale=1.25]
  % diene s-cis: C1 (0,1.6) C2 (0.6,2.4) C3 (1.6,2.4) C4 (2.2,1.6); dienophile below
  \draw[thick] (0,1.6) -- (0.6,2.4) -- (1.6,2.4) -- (2.2,1.6);
  \draw[thick] (0.2,1.62) -- (0.65,2.22);
  \draw[thick] (2.0,1.62) -- (1.55,2.22);
  \draw[thick] (0.2,0) -- (2.0,0); \draw[thick] (0.35,0.14) -- (1.85,0.14);
  \node[font=\scriptsize, anchor=east] at (-0.05,1.6) {1};
  \node[font=\scriptsize, anchor=south east] at (0.5,2.42) {2};
  \node[font=\scriptsize, anchor=south] at (1.6,2.42) {3};
  \node[font=\scriptsize, anchor=west] at (2.25,1.6) {4};
  % arrows: dienophile pi -> C1..C(dienophile); C1=C2 pi -> C2-C3; C3=C4 pi -> C4-C
  \draw[->, thick, omProp] (1.1,0.22) .. controls (0.9,0.9) and (0.3,0.9) .. (0.12,1.42);
  \draw[->, thick, omProp] (0.45,1.95) .. controls (0.6,2.75) and (1.0,2.8) .. (1.1,2.48);
  \draw[->, thick, omProp] (1.78,1.92) .. controls (2.5,1.6) and (2.4,0.6) .. (2.05,0.2);
  \node at (3.3,1.2) {$\longrightarrow$};
  \begin{scope}[xshift=4.2cm]
    \draw[thick] (0,1.6) -- (0.6,2.4) -- (1.6,2.4) -- (2.2,1.6) -- (2.0,0) -- (0.2,0) -- cycle;
    \draw[thick] (0.72,2.26) -- (1.48,2.26);
  \end{scope}
\end{tikzpicture}

{\small \omterm{def:b2:frontier-orbitals:diels-alder}{تفاعل ديلز--ألدر} بين البوتاديين والإيثين، مكتوبًا بالأسهم المنحنية:
تنتقل ثلاثة أزواج من الإلكترونات معًا، في خطوة دائرية متضافرة؛ فتتكوّن رابطتا $\sigma$
عند طرفي \omterm{def:b2:frontier-orbitals:diels-alder}{الديين}، وتنتقل رابطة $\pi$ إلى وسطه.}
\end{omfigure}

التآثر الغالب هو بين المدار \omterm{def:b2:frontier-orbitals:frontier}{HOMO} \omterm{def:b2:frontier-orbitals:diels-alder}{للديين}، $\psi_2$، والمدار \omterm{def:b2:frontier-orbitals:frontier}{LUMO}
\omterm{def:b2:frontier-orbitals:diels-alder}{للمحب للديين}، $\pi^*$. ففصوصهما عند طرفي \omterm{def:b2:frontier-orbitals:diels-alder}{الديين} وعلى كربوني
\omterm{def:b2:frontier-orbitals:diels-alder}{المحب للديين} متوافقة الإشارات حين يقترب \omterm{def:b2:frontier-orbitals:diels-alder}{المحب للديين} من وجه \omterm{def:b2:frontier-orbitals:diels-alder}{الديين}:
فيمكن أن تتكوّن الرابطتان الجديدتان معًا، على الوجه نفسه من كل شريك.

\begin{omfigure}
\begin{tikzpicture}[font=\small]
  \foreach \x/\c in {0/0.602, 1/0.372, 2/-0.372, 3/-0.602} {
    \pic at (\x,2) {omp={1.5*\c}{90}};
  }
  \draw[black!45] (-0.2,2) -- (3.2,2);
  \node[anchor=west] at (3.6,2) {\foreignlanguage{arabic}{مدار \babelsublr{HOMO} للديين، $\psi_2$}};
  \pic at (0,0) {omp={-0.85}{90}}; \pic at (3,0) {omp={0.85}{90}};
  \draw[black!45] (-0.2,0) -- (3.2,0);
  \node[anchor=west] at (3.6,0) {\foreignlanguage{arabic}{مدار \babelsublr{LUMO} للمحب للديين، $\pi^*$}};
  \draw[omProp, thick, densely dashed] (0,1.45) -- (0,0.6);
  \draw[omProp, thick, densely dashed] (3,1.45) -- (3,0.6);
  \node[font=\scriptsize, omProp, anchor=east] at (-0.15,1.0) {\foreignlanguage{arabic}{توافق في الطور}};
  \node[font=\scriptsize, omProp, anchor=west] at (3.15,1.0) {\foreignlanguage{arabic}{توافق في الطور}};
\end{tikzpicture}

{\small المداران الحدوديان للبوتاديين والإيثين، فصوصهما متناسبة مع
معاملات هوكل. وحين تتقابل، يكون لفصي طرفي \omterm{def:b2:frontier-orbitals:frontier}{HOMO} \omterm{def:b2:frontier-orbitals:diels-alder}{الديين} 
ولفصي \omterm{def:b2:frontier-orbitals:frontier}{LUMO} \omterm{def:b2:frontier-orbitals:diels-alder}{المحب للديين} الإشارتان نفسهما عند الطرفين كليهما: فيمكن أن تتكوّن رابطتا
$\sigma$ معًا، على الوجه نفسه من كل شريك.}
\end{omfigure}

\begin{proposition}[النوعية الفراغية]\label{prop:b2:frontier-orbitals:da-stereospecific}
\omterm{def:b2:frontier-orbitals:diels-alder}{تفاعل ديلز--ألدر} نوعي فراغيًا: البدائل الواقعة في وضع cis على \omterm{def:b2:frontier-orbitals:diels-alder}{المحب للديين} تكون
cis في الناتج، والبدائل trans تبقى trans.
\end{proposition}

\begin{proof}
\omterm{def:b2:frontier-orbitals:concerted}{التفاعل متضافر} والرابطتان الجديدتان تتكوّنان على الوجه نفسه من
\omterm{def:b2:frontier-orbitals:diels-alder}{المحب للديين}: فلا يدور كربوناه أحدهما بالنسبة إلى الآخر أبدًا، وتنتقل المواضع
النسبية لبدائلهما إلى الحلقة.
\end{proof}

\begin{proposition}[الانتقائية الموضعية]\label{prop:b2:frontier-orbitals:da-regio}
مع شريكين غير متناظرين، يصل الناتج الرئيسي ذرة \omterm{def:b2:frontier-orbitals:diels-alder}{الديين} ذات
المعامل الأكبر في مدارها \omterm{def:b2:frontier-orbitals:frontier}{HOMO} بذرة \omterm{def:b2:frontier-orbitals:diels-alder}{المحب للديين} ذات المعامل
الأكبر في مدارها \omterm{def:b2:frontier-orbitals:frontier}{LUMO}.
\end{proposition}

\begin{proof}[تعليل]
في الحالة الانتقالية لا تكون الرابطتان الجديدتان قد تكوّنتا بالقدر نفسه بعد؛ 
والاستقرار $2\beta^2/\Delta\varepsilon$ أكبر حين تتداخل المعاملات الأكبر (إذ تزداد $\beta$
مع جداء معاملي الذرتين المتداخلتين)، فيُبلغ ذلك
التوجه بطاقة أدنى.
\end{proof}

\begin{omfigure}
\begin{tikzpicture}
\begin{axis}[ybar, width=0.47\linewidth, height=4.6cm, ymin=-0.8, ymax=0.8,
  xtick={0,1,2,3,4}, xticklabels={O,C1,C2,C3,C4}, axis x line=bottom, axis y line=left,
  extra y ticks={0}, extra y tick style={grid=major, grid style={black!50}}, enlarge x limits=0.15,
  tick label style={font=\scriptsize}, ylabel={\foreignlanguage{arabic}{المعامل}}, label style={font=\small},
  title={\small \foreignlanguage{arabic}{مدار \babelsublr{HOMO} في 1-ميثوكسي البوتاديين}}, bar width=8pt]
\addplot[fill=omDef!50, draw=omDef] table[x=atom, y=c] {figdata/out/bachelor-2-huckel-c-diene.dat};
\end{axis}
\end{tikzpicture}\hfill
\begin{tikzpicture}
\begin{axis}[ybar, width=0.47\linewidth, height=4.6cm, ymin=-0.8, ymax=0.8,
  xtick={1,2,3,4}, xticklabels={C$\beta$,C$\alpha$,C(O),O}, axis x line=bottom, axis y line=left,
  extra y ticks={0}, extra y tick style={grid=major, grid style={black!50}}, enlarge x limits=0.15,
  tick label style={font=\scriptsize}, ylabel={\foreignlanguage{arabic}{المعامل}}, label style={font=\small},
  title={\small \foreignlanguage{arabic}{مدار \babelsublr{LUMO} في البروبينال}}, bar width=8pt]
\addplot[fill=omThm!45, draw=omThm] table[x=atom, y=c] {figdata/out/bachelor-2-huckel-c-dienophile.dat};
\end{axis}
\end{tikzpicture}

{\small نموذج هوكل بوسائط الذرات غير الكربونية ($\alpha_{\ce{O}} = \alpha + 2\beta$،
$\beta_{\ce{CO}} = 0.8\beta$ لأكسجين الإيثر؛ $\alpha_{\ce{O}} = \alpha + \beta$، $\beta_{\ce{CO}} =
\beta$ لأكسجين الكربونيل؛ وهو نموذج). مدار \omterm{def:b2:frontier-orbitals:frontier}{HOMO} \omterm{def:b2:frontier-orbitals:diels-alder}{للديين} أكبر على C4، الطرف
البعيد عن مجموعة الميثوكسي (0.60 مقابل 0.51)؛ ومدار \omterm{def:b2:frontier-orbitals:frontier}{LUMO} \omterm{def:b2:frontier-orbitals:diels-alder}{للمحب للديين} أكبر على
الكربون $\beta$ (0.66). فيرتبطان معًا: وتنتهي مجموعتا الميثوكسي والألدهيد
على كربونين متجاورين من الحلقة.}
\end{omfigure}

ترفع المجموعة المانحة مدار \omterm{def:b2:frontier-orbitals:frontier}{HOMO} \omterm{def:b2:frontier-orbitals:diels-alder}{للديين} ($m = 0.48$ بدل 0.62) وتخفض المجموعة الساحبة
مدار \omterm{def:b2:frontier-orbitals:frontier}{LUMO} \omterm{def:b2:frontier-orbitals:diels-alder}{للمحب للديين} ($m = -0.35$ بدل $-1$): فتضيق الفجوة، وتتفاعل هذه
الأزواج أسرع بكثير من تفاعل البوتاديين مع الإيثين. والتوليفة النموذجية هي \omterm{def:b2:frontier-orbitals:diels-alder}{ديين} غني بالإلكترونات 
\omterm{def:b2:frontier-orbitals:diels-alder}{ومحب للديين} فقير بالإلكترونات.

\begin{definition}[ناتجا الإضافة endo وexo]\label{def:b2:frontier-orbitals:endo}
حين يُضاف \omterm{def:b2:frontier-orbitals:diels-alder}{ديين} حلقي إلى \omterm{def:b2:frontier-orbitals:diels-alder}{محب للديين} يحمل بديلًا غير مشبع، يكون
\emph{ناتج الإضافة endo}\index{ناتج إضافة endo} هو الذي يتجه فيه ذلك البديل نحو الرابطة المزدوجة
المتكوّنة حديثًا (تحت \omterm{def:b2:frontier-orbitals:diels-alder}{الديين} في الحالة الانتقالية)، و\emph{ناتج
الإضافة exo}\index{ناتج إضافة exo} هو الذي يتجه فيه البديل بعيدًا عنها. وتنص \emph{قاعدة endo}\index{قاعدة endo}
على أن ناتج الإضافة endo يتكوّن أسرع.
\end{definition}

\begin{proposition}[قاعدة endo]\label{prop:b2:frontier-orbitals:endo}
تحت التحكم الحركي، يكون \omterm{def:b2:frontier-orbitals:endo}{ناتج الإضافة endo} هو الناتج الرئيسي، مع أن \omterm{def:b2:frontier-orbitals:endo}{ناتج الإضافة exo}
يكون غالبًا هو الأكثر استقرارًا.
\end{proposition}

\admitted

التفسير المعتاد تداخل ثانوي، في الاقتراب endo، بين
فصوص نظام $\pi$ للبديل وفصوص الذرات المركزية في مدار \omterm{def:b2:frontier-orbitals:frontier}{HOMO}
\omterm{def:b2:frontier-orbitals:diels-alder}{للديين}؛ وهو يخفض الحالة الانتقالية endo دون أن يكوّن رابطة. والقاعدة
تجريبية ولها استثناءات؛ ويعالج مجلد السنة الثالثة هذه التفاعلات بواسطة
تناظر مداراتها.

\begin{omfigure}
\begin{tikzpicture}[font=\small]
  \foreach \xs/\mode in {0/endo, 5/exo} {
    \begin{scope}[xshift=\xs cm]
      % cyclopentadiene seen from the side: C1, C2=C3, C4, CH2 bridge on top
      \draw[thick] (0,2) -- (0.6,2.5) -- (1.8,2.5) -- (2.4,2);
      \draw[thick] (0.75,2.38) -- (1.65,2.38);
      \draw[thick] (0,2) -- (1.2,3.2) -- (2.4,2);
      \node[font=\scriptsize, anchor=south] at (1.2,3.2) {\ce{CH2}};
      % dienophile
      \draw[thick] (0.2,0) -- (2.2,0); \draw[thick] (0.35,0.1) -- (2.05,0.1);
      \draw[omProp, densely dashed, thick] (0,1.9) -- (0.2,0.15);
      \draw[omProp, densely dashed, thick] (2.4,1.9) -- (2.2,0.15);
      \node[anchor=north] at (1.2,-1.3) {\mode};
    \end{scope}
  }
  % endo: carbonyls tucked under the diene
  \draw[thick] (0.2,0) -- (0.75,0.95) (2.2,0) -- (1.65,0.95);
  \node[font=\scriptsize, fill=white, inner sep=1pt] at (0.75,1.1) {C=O};
  \node[font=\scriptsize, fill=white, inner sep=1pt] at (1.65,1.1) {C=O};
  \draw[omThm, densely dotted, thick] (0.75,1.25) -- (0.65,2.35);
  \draw[omThm, densely dotted, thick] (1.65,1.25) -- (1.75,2.35);
  % exo: carbonyls pointing away
  \draw[thick] (5.2,0) -- (4.95,-0.85) (7.2,0) -- (7.45,-0.85);
  \node[font=\scriptsize] at (4.95,-1.0) {C=O};
  \node[font=\scriptsize] at (7.45,-1.0) {C=O};
\end{tikzpicture}

{\small اقترابا endo وexo لأنهيدريد الماليك من البنتاديين الحلقي، رسم تخطيطي.
الروابط قيد التكوّن متقطعة. في الاقتراب endo تقع مجموعتا الكربونيل تحت
\omterm{def:b2:frontier-orbitals:diels-alder}{الديين}، وتتداخل مداراتهما $\pi$ مع كربوناته المركزية (منقّط، تداخل
ثانوي)؛ وفي الاقتراب exo تتجهان بعيدًا.}
\end{omfigure}

\begin{method}[رسم ناتج تفاعل ديلز--ألدر]\label{met:b2:frontier-orbitals:diels-alder}
\begin{enumerate}
  \item ضع \omterm{def:b2:frontier-orbitals:diels-alder}{الديين} في تشكّله s-cis، \omterm{def:b2:frontier-orbitals:diels-alder}{والمحب للديين} مقابل طرفيه.
  \item ارسم رابطتي $\sigma$ الجديدتين من طرفي \omterm{def:b2:frontier-orbitals:diels-alder}{الديين} إلى كربوني \omterm{def:b2:frontier-orbitals:diels-alder}{المحب للديين}،
        ورابطة $\pi$ الجديدة بين C2 وC3 في \omterm{def:b2:frontier-orbitals:diels-alder}{الديين}.
  \item أبقِ بدائل \omterm{def:b2:frontier-orbitals:diels-alder}{المحب للديين} في الجهة نفسها (cis يبقى cis).
  \item مع شريكين غير متناظرين، صِل المعاملات الأكبر (نواتج «أورثو» أو «بارا»
        للديينات المستبدلة في الموضع 1 أو 2).
  \item مع \omterm{def:b2:frontier-orbitals:diels-alder}{ديين} حلقي، ضع البديل غير المشبع في الوضع endo.
\end{enumerate}
\end{method}

\begin{inthelab}[تكسير ثنائي البنتاديين الحلقي]
يُباع البنتاديين الحلقي على هيئة ثنائيه ويُستعاد قبيل الاستعمال: يُسخَّن الثنائي
إلى درجة غليانه (نحو \qty{170}{\degreeCelsius}) في دورق مزوَّد
بعمود تجزئة، حيث ينشطر عائدًا إلى الأحادي (\omterm{def:b2:frontier-orbitals:diels-alder}{تفاعل ديلز--ألدر}
عكسي)، الذي يتقطر عند \qty{41}{\degreeCelsius} ويُجمع في مستقبِل مبرَّد
في الجليد. ويعود الأحادي فيتثنّى في غضون ساعات في درجة حرارة الغرفة فيُستعمل
فورًا. والمركبان كلاهما قابلان للاشتعال وضاران؛ ويجري العمل تحت ساحبة أبخرة.
\end{inthelab}

\begin{safety}{\ghs[5mm]{GHS02}\,\ghs[5mm]{GHS06}\,\ghs[5mm]{GHS08}}
البنتاديين الحلقي وثنائيه شديدا الاشتعال وسامّان بالاستنشاق؛ وأنهيدريد الماليك
أكّال ومحسِّس تنفسي. ويُتداولان تحت ساحبة أبخرة،
بقفازات وواقٍ للعينين، بعيدًا عن اللهب.
% ledger: ghs:cyclopentadiene, ghs:dicyclopentadiene, ghs:maleic-anhydride, bp:cyclopentadiene, bp:dicyclopentadiene
\end{safety}

\section{تمارين}

\begin{exercise}[$\star$]\label{exo:b2:frontier-orbitals:1}
أعطِ المدارين \omterm{def:b2:frontier-orbitals:frontier}{HOMO} \omterm{def:b2:frontier-orbitals:frontier}{وLUMO} (الطاقات بوحدات هوكل) للإيثين، ولأنيون الأليل
وللبوتاديين.
\end{exercise}

\begin{exercise}[$\star$]\label{exo:b2:frontier-orbitals:2}
صنّف إلى قاسٍ أو لين: \ce{OH-}، \ce{I-}، \ce{H+}، ثيولات \ce{RS-}، كربون
اليودوميثان، \ce{F-}.
\end{exercise}

\begin{exercise}[$\star$]\label{exo:b2:frontier-orbitals:3}
أيّ هذه الديينات يمكن أن يشارك في \omterm{def:b2:frontier-orbitals:diels-alder}{تفاعل ديلز--ألدر}: البوتا-1،3-ديين،
والبنتاديين الحلقي، والبنتا-1،4-ديين، و(2\textit{Z}،4\textit{Z})-الهكسا-2،4-ديين في هيئته
s-cis؟
\end{exercise}

\begin{exercise}[$\star$]\label{exo:b2:frontier-orbitals:4}
ارسم ناتج البوتاديين مع البروبيننتريل \ce{CH2=CH-C#N}.
\end{exercise}

\begin{exercise}[$\star\star$]\label{exo:b2:frontier-orbitals:5}
مداران عند $-10$ و\qty{-2}{eV} يتآثران مع $\beta = \qty{-1.0}{eV}$ (معطيات التمرين).
احسب الاستقرار الاضطرابي لإلكترونين، ثم الاستقرار المضبوط.
\end{exercise}

\begin{exercise}[$\star\star$]\label{exo:b2:frontier-orbitals:6}
يتفاعل أيون السيانيد مع اليودوميثان فيعطي الإيثاننتريل \ce{CH3-C#N}، لكنه مع
كاتيون كربوني في مذيب شديد التأيين يعطي شيئًا من الإيزونتريل \ce{R-N#C}. فسّر.
\end{exercise}

\begin{exercise}[$\star\star$]\label{exo:b2:frontier-orbitals:7}
مستعملًا معاملات الفصل، تنبأ بالناتج الرئيسي لتفاعل 1-ميثوكسي البوتاديين
مع البروبينال.
\end{exercise}

\begin{exercise}[$\star\star$]\label{exo:b2:frontier-orbitals:8}
ارسم \omterm{def:b2:frontier-orbitals:endo}{ناتج الإضافة endo} للبنتاديين الحلقي مع بروبينوات الميثيل.
\end{exercise}

\begin{exercise}[$\star\star$]\label{exo:b2:frontier-orbitals:9}
لماذا يتفاعل أنهيدريد الماليك مع البنتاديين الحلقي في درجة حرارة الغرفة، بينما يحتاج الإيثين
إلى التسخين تحت الضغط؟
\end{exercise}

\begin{exercise}[$\star\star\star$]\label{exo:b2:frontier-orbitals:10}
يتفاعل البنتاديين الحلقي مع مالييات ثنائي الميثيل (ثنائي إستر cis) ومع فومارات ثنائي الميثيل
(ثنائي إستر trans). ارسم الناتجين وقل كم متماكبًا فراغيًا يعطي كل منهما.
\end{exercise}

\begin{exercise}[$\star\star\star$]\label{exo:b2:frontier-orbitals:11}
\omterm{def:b2:frontier-orbitals:diels-alder}{لتفاعل ديلز--ألدر} $\Delta_r H^\circ < 0$ و$\Delta_r S^\circ < 0$. فسّر الإشارتين
ولماذا يصير التفاعل العكسي مواتيًا في درجة الحرارة المرتفعة.
\end{exercise}

\begin{exercise}[$\star\star\star$]\label{exo:b2:frontier-orbitals:12}
في الهيدرازين \ce{H2N-NH2}، تحمل كل ذرة نيتروجين زوجًا حرًا. مستعملًا تآثر
الإلكترونات الأربعة، فسّر لماذا يتجنب الجزيء التشكّل الذي يكون فيه الزوجان
الحران متوازيين.
\end{exercise}

\section{مسألة: تكسير البنتاديين الحلقي}

\begin{problem}[{مسألة نهاية الأسبوع --- تثنّي البنتاديين الحلقي بوصفه
تفاعل ديلز--ألدر، وترموديناميك التكسير، والمدارات الحدودية
للبنتاديين الحلقي وأنهيدريد الماليك، والكيمياء الفراغية لناتج إضافتهما}]
\label{pb:b2:frontier-orbitals:1}
المعطيات: درجات الغليان، البنتاديين الحلقي \qty{41}{\degreeCelsius}، وثنائي البنتاديين الحلقي نحو
\qty{170}{\degreeCelsius}. طاقات نموذج هوكل (وهو نموذج): البنتاديين الحلقي، معامَلًا على أنه
وحدة بوتاديين، \omterm{def:b2:frontier-orbitals:frontier}{HOMO} عند $\alpha + 0.62\beta$ \omterm{def:b2:frontier-orbitals:frontier}{وLUMO} عند $\alpha - 0.62\beta$؛ وأنهيدريد الماليك، \omterm{def:b2:frontier-orbitals:frontier}{LUMO}
عند نحو $\alpha - 0.35\beta$.
% ledger: bp:cyclopentadiene, bp:dicyclopentadiene

\medskip
\noindent\textbf{الجزء الأول --- الثنائي.}
\begin{enumerate}
  \item ارسم البنتاديين الحلقي وبيّن أن وحدته الديينية مثبتة في الهيئة s-cis.
  \item في التثنّي، يؤدي أحد الجزيئين دور \omterm{def:b2:frontier-orbitals:diels-alder}{الديين}. فما دور الآخر؟
  \item ارسم الثنائي ورقّم الروابط الجديدة.
  \item أيّ ناتجي الإضافة، endo أم exo، يتكوّن في درجة حرارة الغرفة، ولماذا؟
  \item هل التثنّي نوعي فراغيًا؟ فسّر بالآلية.
  \item لماذا لا يحتاج التثنّي إلى حفاز؟
\end{enumerate}

\medskip
\noindent\textbf{الجزء الثاني --- ترموديناميك التكسير.}
\begin{enumerate}[resume]
  \item أعطِ إشارة $\Delta_r H^\circ$ للتثنّي، انطلاقًا من الروابط المتكوّنة 
        والمنكسرة.
  \item أعطِ إشارة $\Delta_r S^\circ$.
  \item اكتب $\Delta_r G^\circ(T)$ وبيّن أنها تغيّر إشارتها عند درجة حرارة $T_i$.
  \item لماذا يكون الثنائي مستقرًا في درجة حرارة الغرفة، والأحادي مفضّلًا عند التسخين؟
  \item في جهاز التكسير، لماذا يدفع تقطيرُ الأحادي حال تكوّنه
        التفاعلَ إلى الأمام؟
  \item لماذا يجب حفظ الأحادي باردًا واستعماله بسرعة؟
\end{enumerate}

\medskip
\noindent\textbf{الجزء الثالث --- \omterm{def:b2:frontier-orbitals:frontier}{المدارات الحدودية}.}
\begin{enumerate}[resume]
  \item في تفاعل البنتاديين الحلقي مع نفسه، احسب الفجوة HOMO--LUMO بوحدة
        $|\beta|$.
  \item في تفاعل البنتاديين الحلقي مع أنهيدريد الماليك، احسب الفجوة بين مدار \omterm{def:b2:frontier-orbitals:frontier}{HOMO}
        \omterm{def:b2:frontier-orbitals:diels-alder}{للديين} ومدار \omterm{def:b2:frontier-orbitals:frontier}{LUMO} للأنهيدريد.
  \item أي الزوجين يتفاعل أسرع؟ فسّر.
  \item لماذا تخفض مجموعتا الكربونيل في أنهيدريد الماليك مداره \omterm{def:b2:frontier-orbitals:frontier}{LUMO}؟
  \item أي مدار في أنهيدريد الماليك يعطي التداخل الثانوي في الاقتراب
        endo؟
\end{enumerate}

\medskip
\noindent\textbf{الجزء الرابع --- ناتج الإضافة مع أنهيدريد الماليك.}
\begin{enumerate}[resume]
  \item ارسم \omterm{def:b2:frontier-orbitals:endo}{ناتج الإضافة endo}.
  \item علّم مراكزه الفراغية.
  \item هل ناتج الإضافة كيرالي؟ (ابحث عن مستوي مرآة.)
  \item لماذا ينتهي هيدروجينا \omterm{def:b2:frontier-orbitals:diels-alder}{المحب للديين} السابق في الجهة نفسها من
        الحلقة؟
  \item كيف سيبدو \omterm{def:b2:frontier-orbitals:endo}{ناتج الإضافة exo}؟
  \item هل يمكن أن يتحول ناتجا الإضافة endo وexo أحدهما إلى الآخر في درجة حرارة الغرفة؟ ولماذا؟
  \item كيف سيتفاعل الأنهيدريد مع الماء، وما الناتج الذي سيتكوّن؟
  \item اذكر عدد المراكز الفراغية في \omterm{def:b2:frontier-orbitals:endo}{ناتج الإضافة endo} للبنتاديين الحلقي 
        وأنهيدريد الماليك.
\end{enumerate}
\end{problem}
